%! AP Statistics — Unit 1, Lesson 1.6 — Exit Ticket ANSWER KEY
\documentclass[10pt]{article}
\usepackage{apstats-article}
\usepackage{apstats-key}

%=============================================================================
\begin{document}

\pageheader{Unit 1, Lesson 1.6}{Exit Ticket --- Answer Key}
\begin{tabularx}{\linewidth}{X X X}
Name:\;\ans{Key} & Date:\; & Period:\;
\end{tabularx}

\vspace{0.22in}

\begin{tcolorbox}[colback=sky, colframe=navy, boxrule=0.9pt, arc=2mm,
    left=4mm, right=4mm, top=3mm, bottom=3mm]
\small Commute times (minutes) for 20 workers: skewed right, center $\approx$ 22
min, range 8--88 min, two outliers above 75 min.
\end{tcolorbox}

\vspace{0.2in}

\begin{enumerate}[leftmargin=1.5em, itemsep=10pt]

  \item \ans{The distribution of commute times for these 20 workers is skewed
        right, with most workers commuting between 10 and 30 minutes and a long
        tail toward higher commute times. [S] \quad There appear to be two
        potential outliers --- workers with commutes over 75 minutes --- which
        stand out well above the main cluster of the data. [O] \quad The center
        of the distribution is approximately 22 minutes, suggesting that a
        typical worker in this group commutes about 22 minutes. [C] \quad The
        spread is large: commute times range from 8 to 88 minutes, a range of
        80 minutes. [S]}

  \begin{teachernote}
  AP rubric: 1 pt each for (S) skewed right with justification, (O) outliers
  identified with approximate value, (C) center with units and context,
  (S) spread described as a range with units. Deduct for missing context in
  any sentence. Accept minor variation in center/spread estimates.
  \end{teachernote}

  \item \ans{The classmate is incorrect. For a skewed-right distribution, the
        few very large values (the long tail on the right) pull the mean toward
        higher values, so the mean is typically \textit{greater than} the
        median --- not less than. The median is resistant to outliers, so it
        stays near the center of the bulk of the data.}

  \begin{teachernote}
  A common error: students confuse which measure is larger. Skewed right
  $\Rightarrow$ mean $>$ median because the mean is pulled in the direction
  of skew. Skewed left $\Rightarrow$ mean $<$ median.
  \end{teachernote}

\end{enumerate}

\vspace{0.18in}

\begin{tcolorbox}[colback=redbg, colframe=redacc, boxrule=0.8pt, arc=2mm,
    left=4mm, right=4mm, top=3mm, bottom=3mm]
\small\textbf{AP-Style Practice --- Key}

\vspace{8pt}
\ans{The distribution of ages of these 15 dogs at the animal shelter is
bimodal, with one peak near 2--3 years and a second smaller peak near 8--9
years, suggesting two distinct groups of dogs in the shelter. [S] \quad There
do not appear to be any outliers, as all ages fall within the overall pattern
of the distribution. [O] \quad The center of the distribution is approximately
4--5 years, though the bimodal shape makes a single center value less
meaningful. [C] \quad The spread of the distribution ranges from 1 to 10 years,
a range of 9 years. [S]}

\begin{teachernote}
Key teaching moment: when a distribution is bimodal, the center is less
meaningful as a single summary. Strong responses will acknowledge this.
Award full credit for SOCS with context; bonus credit for noting the
reduced usefulness of a single center value.
\end{teachernote}
\end{tcolorbox}

\end{document}
